% §1 引言

高级持续性威胁（APT）不是孤立的事件——它们是因果链。攻击者利用Web服务器的漏洞，派生出Shell，下载载荷，读取凭证，将数据外泄到外部主机。每一步\emph{导致}下一步。对于调查告警的安全分析员来说，知道\emph{哪条因果相连的事件序列}构成了攻击，比知道哪个15分钟时间窗口或哪个单独文件可疑更具可操作性。

溯源图由内核级审计日志构建，天然捕获了这些因果依赖关系。节点代表系统实体（进程、文件、套接字），有向边代表带有精确时间戳的因果系统事件（进程派生、文件写入、网络发送）。由于内核调解所有系统交互，溯源图提供了完整且防篡改的因果系统活动记录。

尽管运行在显式编码因果的数据结构上，基于溯源的入侵检测系统（PIDS）尚未实现因果链检测的潜力：

\begin{enumerate}
\item \textbf{粗粒度检测。} 最先进的PIDS在时间窗口~\cite{cheng2024kairos}、单个实体~\cite{jia2024magic,wang2022threatrace}、图快照~\cite{yang2023prographer}或边级别~\cite{edgetrace2025}进行检测。没有任何方法识别哪条具体的因果相连的事件序列构成了攻击。
\item \textbf{启发式链提取。} 最近的Transformer-on-provenance方法使用固定时间窗口~\cite{eagleeye2024}、随机游走~\cite{sentient2026}或时序子图~\cite{getaid2025}来分割溯源图。它们都不尊重因果边方向——固定窗口混合了因果无关的事件，随机游走可能沿逆向遍历边，时序子图使用时间而非因果作为组织原则。
\item \textbf{后处理重建。} 产生攻击路径的方法是在节点/边级检测之后作为后处理步骤进行的~\cite{slot2025,cage2025}。链质量从根本上受限于上游检测的召回率——漏掉一个关键节点就会导致链断裂。
\end{enumerate}

一个更深层的问题潜藏于所有这些局限之下：\textbf{没有任何已有方法在提取链之前定义了什么是“因果上好”的链。}已有方法将链提取视为一个分割问题——如何将大图切成小片段——并依赖下游模型性能间接验证分割质量。

\textbf{我们的方法。} 我们将链提取重新定义为\emph{度量}问题。我们提出一个模型无关的度量，直接量化从溯源图中提取的任意事件链的因果结构完整性，以及一个仅从良性数据中发现该度量参数的校准协议。

具体而言，我们做出了以下贡献：

\begin{enumerate}
\item \textbf{因果相干性度量（$\branchCoherence$）。} 一个模型无关的度量，根植于操作系统原理：进程是独占的能力（内核保证进程身份），而文件和套接字是共享资源（受第三方修改影响）。该度量采用三分法的逐对相干性函数和惩罚高分支良性进程的结构稀疏先验。
\item \textbf{无监督校准协议。} 基于肘部检测的方法，仅从良性溯源数据中发现操作系统的自然因果视野，零攻击标签、零人工参数调优。该协议具有跨平台性：自动适配不同操作系统之间审计日志粒度的差异。
\item \textbf{端到端检测系统。} 流水线组合了基于TGN的种子识别（复用KAIROS）、以$\branchCoherence$为指导的因果链提取器、以及通过自回归预测事件过渡来区分攻击链与良性链的自监督Transformer。Transformer仅在良性数据上训练。
\item \textbf{实证验证。} 在DARPA透明计算基准上，因果链达到[X]\%的链级F1，超过基于窗口、基于随机游走、基于子图和基于规则的链提取方法超过10个百分点的绝对F1，同时相比KAIROS减少超过50倍的调查成本。校准协议在不同数据集（E3和E5）上发现不同的因果视野，展示了跨平台自动适配能力。
\end{enumerate}

本文其余部分组织如下。\S\ref{sec:background}提供溯源图背景并分析已有链提取策略的局限性。\S\ref{sec:related}将我们的工作置于更广泛的溯源入侵检测文献中。\S\ref{sec:design}展示系统设计，详细讨论因果链提取器、相干性度量和自回归Transformer。\S\ref{sec:metric}在独立于任何模型训练的情况下验证相干性度量。\S\ref{sec:evaluation}呈现端到端检测结果。\S\ref{sec:discussion}讨论部署考量和局限性。\S\ref{sec:conclusion}总结。
